\documentclass[aspectratio=169]{beamer}
\usepackage{graphicx}
\usepackage{natbib}
\usepackage{caption}
\usepackage{pdflscape}
\usepackage{color}
\usepackage{soul}
\usepackage{amsmath,amssymb,graphicx,geometry,gensymb,pifont}
\usepackage[section]{placeins}
\usepackage[thinlines]{easytable}
\usepackage{fancyhdr}
\usepackage{float}
\usepackage{ragged2e}
\usepackage[listings,theorems]{tcolorbox}
\usepackage{subcaption}
\usepackage{animate}

\usepackage[ruled,vlined,linesnumbered,noend]{algorithm2e}
\SetAlFnt{\footnotesize}
\setlength{\algomargin}{0.6em}
\DontPrintSemicolon
\SetKwInput{KwIn}{Input}
\SetKwInput{KwOut}{Output}

\usepackage{tabularray}
\usepackage{tabularx}
\usepackage{booktabs}      % pretty rules
\usepackage{tabularx}      % automatic-width column(s)
\usepackage{makecell}      % easy line breaks inside cells
\renewcommand\theadfont{\normalsize\bfseries}
\renewcommand\cellalign{tl}  % left-align makecell content


\captionsetup[figure]{font=tiny}



\definecolor{maroon1}{rgb}{0.478, 0, 0.235}
\definecolor{dimgray}{rgb}{0.369, 0.416, 0.443}
\definecolor{lightgray}{rgb}{0.94,0.94,0.97}
\definecolor{lightgray2}{rgb}{0.5,0.5,0.5}
\definecolor{darkgreen}{rgb}{0, 0.5, 0}
\definecolor{lightpink}{rgb}{1, 0.93, 0.96}



\usetheme{Copenhagen}
\usecolortheme[named=maroon1]{structure}

\setbeamerfont{caption}{size=\scriptsize}

\setbeamertemplate{navigation symbols}{}

% \expandafter\def\expandafter\insertshorttitle\expandafter{%
% \insertshorttitle\hfill%
% \insertframenumber\,/\,\inserttotalframenumber}

\setbeamercolor{block body}{fg=black,bg=lightgray}

\title[March 2026]{{\bf A Practical Reinforcement Learning Control Framework for Complex Process Systems}}
\author[A. H. Hamedi]{\textcolor {black} { \underline{\bf \scriptsize A. H. Hamedi}\\\vspace{0.2cm}{\scriptsize \color{maroon1} Department of Chemical Engineering, McMaster University}\\\vspace{0.1cm}}}


\date[\today]{}

\begin{document}

%%%%%%% One-slide summary %%%%%%%%
\begin{frame}
\frametitle{\centering \scriptsize \textbf{Practical Reinforcement Learning Control Framework for Complex Process Systems (Amir H. Hamedi)}}
\scriptsize

\vspace{-0.5cm}

\begin{columns}[T,onlytextwidth]

%========================
% Left: RL-assisted MPC
%========================
\begin{column}{0.48\textwidth}

\begin{block}{\textbf{RL-assisted MPC} \scriptsize{(ongoing work)}}
\begin{itemize}\justifying
    \item RL improves \textbf{linear MPC online} with multi-agent structure.
\end{itemize}
\end{block}

\vspace{-0.15cm}
\begin{figure}
    \centering
    \includegraphics[width=.7\linewidth,height=0.25\textheight]{Figures/RL Assisted.png}
\end{figure}

\vspace{-0.5cm}
\begin{center}
{\tiny \textbf{Proposed RL-assisted MPC diagram}}
\end{center}

\vspace{-0.3cm}
\begin{figure}
    \centering
    \includegraphics[width=0.78\linewidth,height=0.3\textheight]{Figures/compare_outputs_last_episode_nominal.png}
\end{figure}

\vspace{-0.7cm}
\begin{center}
{\tiny \textbf{Polymer reactor: RL-assisted MPC (blue), MPC (orange)}}
\end{center}

\end{column}

%========================
% Right: Practical RL controller
%========================
\begin{column}{0.49\textwidth}

\begin{block}{\textbf{Practical RL controller}}
\begin{itemize}\justifying
    \item Start from \textbf{offset-free linear MPC}, then improve online for complex process systems.
\end{itemize}
\end{block}

\vspace{-0.2cm}
\begin{figure}
    \centering
    \includegraphics[width=0.78\linewidth,height=0.35\textheight]{Figures/C2_nom_outputs_compare_three.png}
\end{figure}

\vspace{-0.22cm}
\begin{center}
{\tiny \textbf{C$_2$ splitter: Modified RL (blue), Pre-trained RL (orange), MPC (green dotted)}}
\end{center}

\vspace{-0.22cm}
\begin{block}{\textbf{Ongoing work}}
\centering
\scriptsize Stable constrained RL using a Lyapunov contraction idea and an output-constraint safety layer
\end{block}

\end{column}

\end{columns}

\end{frame}

\end{document}